% Independently audited by path_stability and random_hosts.
% Exact certificates: four-regular-five-pairing-dual-checks.json.
\section{The sharp error order for cyclic cuts of size at most five}
\label{sec:five-cut-universal-crossing}

The five-port crossing bound has a short distinct-representative proof.
It uses the protected small-cut theorem and the exact five-cut identity.
No terminal-state profile classification or aggregate estimate is needed.

\begin{lemma}[Distinct representatives for five pairings]
\label{lem:five-pairing-distinct-representatives}
For each four-element subset $S\subset[5]$, choose a pairing $\alpha_S$ of $S$.
Regard $\alpha_S$ as two disjoint edges of $K_5$.
One can select an edge $e_S\in\alpha_S$ for each $S$ so that the five selected edges are distinct.
\end{lemma}

\begin{proof}
Form an auxiliary graph whose vertices are the ten edges of $K_5$.
Two auxiliary vertices are adjacent when their corresponding edges are disjoint.
This graph is the Petersen disjointness graph $KG(5,2)$.
The five chosen pairings give five distinct auxiliary edges.
Indeed, the union of the two endpoints determines the four-set $S$.

The auxiliary graph has no triangle, since three disjoint two-sets require six points.
It has no four-cycle either.
Two distinct two-sets have a union of size at least three, so there is at most one two-set disjoint from both.
Thus two distinct auxiliary vertices cannot have two common neighbors.

A five-edge subgraph of this auxiliary graph is therefore a forest, unless all five edges form one five-cycle.
For each tree component, choose a root and assign every edge to its endpoint farther from the root.
For a five-cycle, orient the cycle and assign each edge to its head.
These assignments use distinct auxiliary vertices.
They give the required distinct representative edges of $K_5$.
\end{proof}

\begin{lemma}[A universal five-port functional]
\label{lem:universal-five-port-functional}
Let $a_e\ge0$ be indexed by the ten pairs of $[5]$.
For each four-set $S$, let $u_{S,\alpha}\ge0$ be indexed by its three pairings, and put $U_S=\sum_\alpha u_{S,\alpha}$.
Suppose
\[
 a_e\le\tfrac12+\eta_1,
 \qquad
 a_e+U_S-u_{S,\alpha}\le1+\eta_2
 \quad(e\in\alpha),
\]
where $\eta_1,\eta_2\ge0$.
Then
\[
 \sum_e a_e+\sum_S\max_\alpha(U_S-u_{S,\alpha})
 \le\frac{15}{2}+5\eta_1+5\eta_2.
\]
\end{lemma}

\begin{proof}
For each $S$, choose a maximizing pairing $\alpha_S$.
Apply Lemma~\ref{lem:five-pairing-distinct-representatives} and choose distinct $e_S\in\alpha_S$.
Summing the five corresponding augmented inequalities gives
\[
 \sum_S\bigl(U_S-u_{S,\alpha_S}\bigr)
 +\sum_S a_{e_S}\le5+5\eta_2.
\]
Exactly five path coordinates remain.
Their total is at most $5/2+5\eta_1$.
Adding them proves the stated inequality.
\end{proof}

\begin{proposition}[A uniform deficit across a five-edge cut]
\label{prop:five-cut-universal-deficit}
Let $G$ have no cyclic edge cut of size at most four.
Suppose $G$ is obtained by joining the ports of normalized five-poles $X,Y$ of ranks $h,k$.
If $k\to\infty$, then their normalized crossing contribution is at most
\[
 \frac{15}{2}+o(1).
\]
The estimate is uniform in $X$ and its rank $h$.
\end{proposition}

\begin{proof}
Normalize the states of $Y$ by $2^k$ and write them as $a_e,u_{S,\alpha},U_S$.
Corollaries~\ref{cor:cyclic-five-path-density} and~\ref{cor:five-pole-unconditional-added-edge} give the hypotheses of Lemma~\ref{lem:universal-five-port-functional}, with errors tending to zero as $k\to\infty$.

Normalize the path coordinates of $X$ by $2^h$ and write them as $b_e$.
They all lie in $[0,1]$.
Write its normalized pairing coordinates as $v_{S,\alpha}\ge0$.
For each $S$, their sum is at most one.
The exact crossing contribution is
\[
 \mathcal B(X,Y)
 =\sum_e b_ea_e
  +\sum_S\sum_\alpha v_{S,\alpha}
       (U_S-u_{S,\alpha}).
\]
Every coefficient $U_S-u_{S,\alpha}$ is nonnegative.
The coordinate bounds on $X$ therefore give
\[
 \mathcal B(X,Y)
 \le\sum_ea_e+\sum_S\max_\alpha(U_S-u_{S,\alpha})
 \le\frac{15}{2}+o(1).
\]
All errors concern only the rank of $Y$.
This proves the stated uniformity.
\end{proof}

\begin{theorem}[Sharp-order five-cut error]
\label{thm:sharp-order-five-cut-error}
There is an absolute constant $K_5$ such that every connected loopless cubic multigraph of rank $r$ with a cyclic edge cut of size at most five satisfies
\[
 c(G)\le2^{r-1}+K_5r2^{r/2}.
\]
The factor $r$ cannot be replaced by $o(r)$ throughout this class.
Moreover, for all sufficiently large $r$, a graph with a cyclic five-edge cut and no smaller cyclic cut satisfies
\[
 c(G)<2^{r-1}.
\]
\end{theorem}

\begin{proof}
If $G$ has a cyclic cut of size at most four, apply Theorem~\ref{thm:sharp-order-four-cut-error}.
Suppose no such cut exists, and choose a cyclic five-edge cut.

Both shores are connected.
Indeed, any component containing a cycle must have at least five boundary edges, since its complement contains the opposite shore's cycle.
It uses all five cut edges.
Any other component on its shore would have no edge to its complement, contrary to the connectedness of $G$.

No vertex has two or more incidences in the cut.
Such a vertex has internal degree at most one.
Moving it to the opposite shore preserves a cycle on both sides and reduces the boundary to at most four.
This would contradict the hypothesis.
Thus the two shores are normalized five-poles $X,Y$.
Write their ranks as $h,k\ge1$ and assume $h\le k$.
Then
\[
 r=h+k+4,\qquad k\ge(r-4)/2.
\]

The uniform crossing bound gives
\[
 \frac{c(G)}{2^{h+k}}
 \le\frac{15}{2}+o(1)
    +\frac{c(X)+c(Y)}{2^{h+k}}.
\]
The affine-space bound gives $c(X)\le2^h-1$.
The asymptotic cycle bound gives $c(Y)\le(1/2+o(1))2^k$.
Since $h\ge1$, their normalized total is at most
\[
 2^{-k}+(1/2+o(1))2^{-h}
 \le\frac14+o(1).
\]
Consequently,
\[
 \frac{c(G)}{2^{h+k}}\le\frac{31}{4}+o(1)<8
\]
for all sufficiently large $r$.
This proves the strict one-half bound in the absence of smaller cyclic cuts.

The remaining bounded ranks are absorbed into $K_5$, using $c(G)\le2^r-1$.
The cyclic-square examples already have a cyclic four-edge cut and excess $\Theta(r2^{r/2})$.
They also belong to the present class, so the polynomial factor is sharp.
\end{proof}

The unresolved sharp-order case is now confined to cubic kernels with no cyclic edge cut of size at most five.
The unrestricted polynomial-error theorem applies to those kernels, but does not give the half-exponent error there.
